\documentclass[12pt, a4paper]{article}

\usepackage[utf8]{inputenc}
\usepackage[T1]{fontenc}
\usepackage[french]{babel}
\usepackage{geometry}
\geometry{top=2.5cm, bottom=2.5cm, left=3cm, right=3cm}

\title{\textbf{L'impact des réseaux sociaux sur la réussite académique}}
\author{Département des Sciences de l'Éducation}
\date{Mars 2026}

\begin{document}

\maketitle

\section{Introduction}

L'essor des réseaux sociaux a profondément modifié le quotidien des étudiants. Facebook, Instagram, TikTok ou YouTube sont devenus des espaces de vie à part entière, consultés plusieurs heures par jour. Cette omniprésence soulève une question centrale : ces plateformes nuisent-elles à la réussite académique, ou peuvent-elles, dans certains cas, la soutenir ?

\section{Effets négatifs sur les apprentissages}

Les effets néfastes les mieux documentés concernent l'attention. Les notifications fragmentent les plages de travail et il faut en moyenne vingt minutes pour retrouver un niveau de concentration équivalent après une interruption numérique. L'usage passif des plateformes favorise également la procrastination : la gratification immédiate offerte par un fil d'actualité entre en compétition directe avec les activités d'étude, dont la récompense est différée. Une méta-analyse portant sur cinquante-neuf études estime qu'une heure supplémentaire par jour sur les réseaux sociaux est associée à une baisse d'environ 0,4 point de moyenne générale.

À cela s'ajoutent les troubles du sommeil causés par l'utilisation des écrans en soirée, qui dégradent la mémorisation et la pensée critique, ainsi que les phénomènes de comparaison sociale — particulièrement intenses sur Instagram — qui peuvent générer anxiété et démotivation.

\section{Usages bénéfiques}

L'image est cependant plus nuancée. Des groupes de travail sur WhatsApp ou Discord facilitent l'entraide et le partage de ressources entre pairs. YouTube s'est imposé comme un outil pédagogique majeur : en 2024, près de trois quarts des étudiants français déclaraient l'utiliser régulièrement à des fins d'apprentissage. Pour les étudiants isolés ou en formation à distance, les réseaux sociaux maintiennent un sentiment d'appartenance communautaire positivement corrélé à la persévérance académique.

\section{Discussion}

L'impact des réseaux sociaux n'est pas uniforme : il dépend du profil de l'utilisateur, du type de plateforme et du contexte d'usage. Les études distinguent généralement trois profils. L'utilisateur non régulé, dont l'usage est passif et subi, enregistre les performances les plus faibles. L'utilisateur partiellement régulé alterne usages utiles et distracteurs. L'utilisateur régulé, dont l'usage est intentionnel et délimité, peut tirer un bénéfice net des outils numériques. Ce dernier profil souligne l'importance de l'autonomie et de la conscience métacognitive dans la relation aux plateformes.

\section{Conclusion}

Les réseaux sociaux ne sont ni des obstacles absolus ni des catalyseurs automatiques de la réussite académique. Leur effet dépend avant tout de la manière dont ils sont utilisés. Former les étudiants à un usage conscient et stratégique du numérique constitue aujourd'hui l'un des enjeux éducatifs les plus pressants.

\end{document}